\section{Germanic Rome}

This translation is from the chapter on the theories of Count de Gobineau from Il mito del sangue, (The Myth of Blood). Although there is much to be said about this passage, I'll leave it to the readers to draw out some conclusions. In some respects Gobineau agrees with Evola: Rome declined and the Middle Ages was the recovery of Tradition by the joint efforts of the Romans and Nordics. After that period, the West declined, paradoxically a victim of its own success. For Gobineau, the spiritual aspect is irrelevant as only the race of the body accounts for everything. Hence, the choice between paganism and Christianity does not matter, at least insofar as the latter was purified of Semitic elements under the influence of the Germanics. What is of interest is Evola's concluding remarks; it probably explains why he had to reject zoological racism. Since this chapter is of significant interest, I hope to post more of it in the next week or two.


\begin{quotationx}
Another product of the mixing of blood for \textbf{Count de Gobineau} would be the feeling for one's homeland and authority, that would arise from the union of Aryans and Semites, from a semitic mitigation of the Aryan predilection for isolation, independence and personality. We will see this theme taken up again by some extremist racialist writers in the sense of relating every form of sovereignty and statolatry comprising the ethnic-national elements to something ``semitic''.

Moreover, the expression ``Semitic Rome'' comes up in de Gobineau to designate the imperial period of that civilization; that, ``not in the sense that it indicates a human type identical to that which results from the Ancient Chaldean and Hamitic combinations: but in the sense that ``of the multitudes scattered with Rome's fortunes over all the countries subject to the Caesars, the greater part was more or less stained with black blood and represented thus a combination not equivalent but analogous to the semitic mixing.''

The predominant ``black'' qualities, well containted in certain limits and compensated by means of some white qualities were, for de Gobineau, essential factors in the development of imperial Rome. In more than one point, the position taken by de Gobineau in regards to Christianity seems negative: too much to listen to again, this popular belief, of ``a religion of slaves, humiliating because pacifistic and egalitarian, and, in a word, unworthy of the races that still conserves any scintilla whatsoever of the Aryan flame.'' In any case, Christianity for him was purified little by little from semitic elements and the Greek made itself Roman (Catholicism) and then, from the Romans, it became Germanic.

For de Gobineau, the Germans and the other Nordic stocks of the invasions appear as the ancient races of pure Aryan blood. But, attracted by the mirage of the Roman symbol, it could not escape the destiny of being dissolved in the potent detritus of the mixed races of Rome, among which their energy and their blood had to decline. This assimilation therefore was not so rapid to drag society to the semitic point of departure typical of the low empire: initially, the Germanic elements were able to be absorbed, but not up to such a sign.

It is thus that the civilisation of ``Germanic Rome'' arises, that is, the Medieval civilisation. Every normal society, for de Gobineau, is founded on three originary classes or castes, corresponding to distinct ethnic strata: ``The nobility, an image more or less like the glorious race; the bourgeoisie, composed of metics similar to the great race; the people, a servile social class belonging to an inferior human type; black in the South, Finnish in the North.'' The Middle Ages still knew such a partition. But it revealed itself always more deprived of its racial base and therefore of its strength. So this hierarchical picture had to unravel little by little while the last veins of pure Aryan blood were snuffed out and scattered. This moves us towards the modern ``repugnant atmosphere of the democratic muck.''

The conclusion of de Gobineau's viewpoint, which is found in his principle work, the famous Essay on the Inequality of the Human Races, that was published between 1853 and 1855, is pessimistic. The dominant impulse of the white race, throwing itself over the entire earth, broke the last ethnic barriers, created a world in which distances no longer exist and where rapprochement, aggregation, and confusion of types are fatal and as rapid as ever. ``Pure Aryans are no longer found''. It is an inexorable law that any group that has the potency of civilisation attracts other races, extends itself, brings itself geographically further and further, dissipates itself, degrades itself. De Gobineau, at the end of his book, says that the history of the world turns ostensibly, through such a path, toward that ``supreme unity'' that, otherwise, he already declared that it was true only for the metics without race.
\end{quotationx}



\flrightit{Posted on 2011-10-17 by Aeneas}

\begin{center}* * *\end{center}

\begin{footnotesize}
\begin{sffamily}
\texttt{logres on 2011-10-18 at 08:44 said:}

Given the nature of genetics (on Gobineau's own assumptions) couldn't you always have a hidden persistence, then resurgence, of certain characteristics or traits?

\hfill

\texttt{Cologero on 2011-10-18 at 23:46 said:}

Gobineau had no understanding of genetics, nor did any of the authors critiqued in the Myth of Blood. As for those traits, the question is whether they are simply genetic traits, or, more likely, related to qualities of the soul or spirit. What is really of interest here is Gobineau's description of the stages in the Christianisation of Europe.

Beginning with the decline of Imperial Rome, ``Semitic Rome'', Gobineau proposes the following stages:

\begin{enumerate}

\def\labelenumi{\Roman{enumi}.}
\item Semitic (primitive Christianity)

\item Greek (Orthodox)

\item Roman (Catholicism)

\item Nordic-Roman (Medievalism)

\item Nordic (Lutheranism)
\end{enumerate}


This is how the later Nordic racialists, for example, Herman Wirth and H S Chamberlain, regarded it; reading Internet forums shows the persistance of this viewpont. These stages correspond remarkably to their views on race. They despise primitive Christianity as Semitic and the product of the lowest classes of Rome. Slavs, for the most part Orthodox, were not considered European.

They are more ambiguous about the Catholic countries of Southern Europe, not wholly approving nor disapproving.

However, in Gobineau's view, the Reformation was not a progressive development, but rather a symptom of racial decline. The Reformation undid the Nordic-Roman creation of a Traditional civilization. Paradoxically and incoherently, the Nordicists are attached to the Reformation, while we would make these points.

\begin{itemize}
\item Luther rejected the Medieval Church for its pagan elements.

\item Luther claimed to be returning to primitive Christianity

\item Luther rejected the Christian OT (Septuagint) in favour of the Jewish Masoretic text
\end{itemize}


This is what perplexed Evola. Instead of going in the direction of a continued repaganisation of Christendom, Luther set in motion the Semiticization of Europe. This movement became progressively stronger in Calvin, the Puritans, and especially now in the zionism of the neo-Christians today. In the other direction, as the force of the Christian element weakened in the North, there developed a philosophical naturalism that was just as opposed to the supernatural and traditional values of Christendom.

This has to be seen as a decline in the race of the soul of Northern Europeans, not in the race of the body. First of all, the Reformation took root in the North, where the effects of race mixing was presumably much less than in the South. I am not convinced by Gobineau's assertion that the North declined from interbreeding with the Finnish element.

Yet, this is what most WNs mean by saving Western civilisation: the re-semiticized Europe both spiritually and philosophically. Gornahoor intends instead, the recovery of Medievalism. The former group has demonstrated little capacity for dialogue on this central point, since it doesn't fit into their preconceptions.

\hfill

\texttt{logres on 2011-10-19 at 00:14 said:}

Perhaps it was simply a case of occult forces playing one against the other, to the maximal detriment of both. A great argument for transcendence of self if there ever was one. Very helpful addendum, thank you. It seems Protestantism cast its lot in with the new money markets and materialism.

\hfill

\texttt{Cologero on 2011-10-19 at 00:54 said:}

It is always a case of occult forces, the third dimension of history, which we have try to identify. The first step is the greater battle, that is, to recognize these forces within oneself. Then the lesser battle is to note the same forces in the external world.

The first is the struggle between cosmos and chaos, or hierarchy vs the undifferentiated mass. This appears as the contest between order and disorder. Insofar as disorder predominates, it indicates a privation, or lack of power and strength of will.

The battle now, however, is between supernaturalism and naturalism. The latter denies the third dimension of occult forces, so the true battle is never even recognized. Instead, there are many pseudo-struggles within naturalism that sap energy and resources. So, yes, it is an ``argument'' for developing the awareness of transcendence.

\hfill

\texttt{Izak on 2011-10-19 at 23:52 said:}

``It seems Protestantism cast its lot in with the new money markets and materialism.''

I'd say this is a solid understatement, if anything. Just look at how reliant Lutheranism was on the printing press, for instance. Without that possibility for mass dissemination of information, you would have none of the protestant insistence on people of lower orders to interpret the word for themselves. They simply wouldn't have had the means.

\hfill

\texttt{Charlotte Cowell on 2011-10-20 at 09:12 said:}

`people of lower orders', who decides which ones fall into THAT category?! (anyone without royal blood, perhaps, would that mean most of us here?). Also remember that people of ahem `higher orders' are more than capable of misinterpretation if their own egos are insufficiently purified or they have personal agendas to pursue. The Holy Spirit speaks to the heart and anyone can listen if they're so inclined. No, I'm not a protestant.

\hfill

\texttt{Graham on 2011-10-20 at 17:24 said:}

Charlotte I believe that's an emotional response to words you don't like. The Catholic point of view holds that the laity should not interpret Scripture for themselves; Scripture is intepreted by Catholic Tradition. Some members of the laity will possess qualifications, no doubt, but that isn't what's at issue when someone speaks against private interpretation.

\hfill

\texttt{Izak on 2011-10-20 at 17:31 said:}

It's a moot point. The lowest order is the third estate, and they didn't interpret the word qua ``The Word'' because they were mostly illiterate and there was no means of mass circulation anyhow. The Catholic church provided images within the church buildings that illustrated the stories, so even the lowest order of all still had to do some degree of interpreting, or ``glossing'' of the images with which they were provided --- and if you go into the ways allegory would be used in story-telling, they still had to do some leg-work there as well. The third estate was challenged intellectually in its religious interpretation, but in a way that suited it better.

But the fact of the matter is that a huge chunk of Protestantism was coming from the fact that the third estate, particularly the mercantile branch, was rising up, the printing press was a new invention, and materialism was beginning to pervade the West. I think Oswald Spengler put it best: while it is true that Protestantism did alleviate the peasant classes from the abuses coming from church authorities (and they were very real), it also put a great deal of weight on their shoulders by setting a mostly unattainable intellectual standard for them to pursue --- most of it dealing with literacy and interpretation. No traditionalist would claim that the word of God should be presented equally to each estate or caste; it's a modern and quasi-democratic idea.

\hfill

\texttt{Charlotte Cowell on 2011-10-21 at 06:03 said:}

Graham, no that wasn't an emotional response although I accept it may have come across in that way. I DO have emotional responses on a regular basis but that doesn't prevent me from rational thinking, I was expressing an actual viewpoint on what Izak rightly describes as an endlessly moot point. Catholic scripture is interpreted by catholic tradition, but the point I made about the `ones in charge', still stands. I personally won't attend any church unless I can doubtlessly trust the priest in terms of sincerity, spiritual integrity and also general intelligence. In other words, unless I can truly accept that they ARE of a higher order. (this does happen of course, but not always). Given a choice I would normally choose to attend a Catholic church as I do in fact believe the catholic church is custodian in no small degree of the Christian ministry. Just to clarify my position, I was converted to Christianity by a Benedictine, I've visited the Vatican twice, own a picture of the Pope and numerous icons of the blessed Virgin and was present at the angelic prayers offered upon the death of John Paul II. No fear of truth here... \hfill

\texttt{Charlotte Cowell on 2011-10-21 at 06:14 said:}

izak, thank you for the further explanation, as it happens I'm kind of in agreement with you and think very valid points are made. Perhaps after all there is some emotion in the response, as Graham suggested, but not to the point of irrationality. I'm frustrated at the deliberate mis-education of people -- the deliberate dumbing down of society -- that is being perpetrated in the world right now. However, at this point in time the church -- Catholic or otherwise -- certainly is NOT the culprit, it's very much a political problem. In the past the church was in control of this situation but that's not the case anymore -- not in most of the western world at any rate. Where the catholic church was at fault in the past (in my view), is where they manipulated the Word in a way that suited political agendas, and used it to instill fear rather than love of God and Truth. There's no point denying this happened as people were forcibly converted -- NOT the way Christ intended. In some cases however I agree that it is much better to withold information than to let people abuse it -- by suppressing grimoires, for instance, but the Word of God is for us all, a small child can understand it perfectly. What happened during the Reformation was a necessary evil in response to sale of indulgences and so on -- we all know what the Medicis and Borgias did, it can't just be accepted as having been acceptable, SOMETHING had to happen to stop it. As with so much of history the pendulum has swung back and forth violently as humanity has struggled to find its balance, but sooner or later it will happen and I do believe that the church is rising to the occasion right now. I just wish there was more stringent criticism of the violent imperial policies of the US, UK, France, etc. What I'm curious about at the moment is what conclusions the Vatican will publically draw on the matter of extra-terrestrial life -- I believe they've been discussing the problem with NASA in the past couple of months?

\hfill

\texttt{Charlotte Cowell on 2011-10-21 at 06:28 said:}

by the way, I would really appreciate some discussion of and insight into what happened during the Vatican II council as the questions to have arisen from this present me with a big stumbling block. If this isn't the right place then perhaps someone knows of where it's being discussed? I've looked into what Malachi Martin has to say as he appears to be very close to the source, what are your thoughts on him?

\hfill

\texttt{Charlotte Cowell on 2011-10-21 at 06:42 said:}

Also (sorry, on a roll now), we shouldn't be afraid to express emotion in relation to a religion that is founded on Love, though certainly you're right in saying that we shouldn't let passion cloud judgement -- the problems of the world make that clear enough. I like the assessment made my Khalil Gibran in the Prophet:

``Your reason and your passion are the rudder and sails of your seafaring soul. If either your sails or your rudder be broken, you can but toss and drift, or else be held at a standstill in mid-seas.

``For reason, ruling alone, is a force confining; and passion, unattended, is a flame that burns to its own destruction.

``Therefore let your soul exalt your reason to the height of passion, that it may sing. And let it direct your passion with reason, that your passion may live through its own daily resurrection, and like the phoenix rise above its own ashes.''

This is probably why your guys considered turning to Islam -- the way of Beauty is more than compelling -- but it's not incompatible with Christianity, there shouldn't be a conflict or dilemma. Moon and Sun, both with one Father and a sole creation.

\hfill

\texttt{Cologero on 2011-10-21 at 08:29 said:}

Hola, Charlotte.

Gornahoor follows Tradition and is opposed to the alleged progressivism of the modern world. Vatican II, with its \emph{aggiornamento} was a deliberate opening up to the modern world. So there is no debate possible between Tradition and Vatican II because the spiritual world views are so incompatible. There are many sites and forums opposed to V II, though, of course, the majority are in favour. Here is one example of the former: \textit{Maurice Pinay}\footnote{\url{https://mauricepinay.blogspot.com/}}, and I can provide more if you want.

\hfill

\texttt{Charlotte Cowell on 2011-10-21 at 13:22 said:}

Hola C!

Thank you for clarifying that, it confirms what I suspected, which is the depth of the breach. I'm not surprised, it is a big problem and as you know I share the concern... .the horror. That said, my attention has been called, by way of reminder, to the spiritual mission of JPII, who it seems descended into hell in order to perform certain tasks of atonement on behalf of the church. We can make of that what we will. Personally speaking I struggle with it all and have to accept there is an impasse. I turn up at my local Church and instead of going inside to join with the service, I sit in the garden and pray with the statue of our Lady, or wipe cobwebs from the crucified Christ. At Easter I took mass there and spent months feeling as if my Spirit might have been manacled into flesh and that the fairies had all died. Every cell in the human body is renewed in around 6 months. I noted on the next visit that they were keeping the wine in the middle of the church, up front so to speak, and took some heart -- it seemed sweeter than before.

Hoo -- Lord of Hosts, yes, but remember that the Christian community is a spiritual one and that the left hand doesn't know what the right is doing. It may not always be readily apparent who's been so appointed and it should be well known to us all that angels can be entertained unawares.

Other than that, do you not agree with Origen's view of universal salvation?

As for caste, we must take care not to equate it with material wealth and be aware that powerful veins can be filled with black dragon blood.

We should have a think about the Dead Sea Scrolls some day, I heard they were owned by the Rockefellers, but it could be another one of those stories eh?

\hfill

\texttt{Perennial on 2011-10-22 at 04:54 said:}

Charlotte, I have noted before elsewhere on here that many of the allegations of curruption in the Church in the past are either exaggerated or pure fantasy. Therefore, justifications against Her, like the protestant ``reformation,'' are generally unjustified. The main example of the selling of indulgences, by the Dominican priest Johann Tezel, was a rare thing, certainly insufficient to start a reformation, which rebelled against both Church and Empire, as people often forget. Many practices, like the selling of indulgences, in fact, had already been condemned. The situation after Vatican II, however, is another matter. There is certainly much do be discussed there, even if Gornahoor is not the place to do it. Perhaps Vatican II and it's fruits will be mentioned on the Right Hand Path blog at some point.

\hfill

\texttt{Cologero on 2011-10-22 at 09:29 said:}

Vatican II has been discussed ad nauseum and those committed to it are adamant. Nevertheless, a return to the \emph{status quo ante} is insufficient ... a better discussion, I think, would be about what would have to happen at Vatican III to bring it into alignment with Tradition.

\hfill

\texttt{Logres on 2011-10-22 at 11:21 said:}

I would love to rejoin the Catholic Church, except, I have no idea what I would be joining, and whether it would be out of the frying pan, and into the fire. The potential for a bad pope to show up and basically make VII irreversible is enormous (from my perspective), while what seems more likely is a slow and incremental giving way to modern pressure. But that's admittedly an outsider-perspective.

\hfill

\texttt{Charlotte Cowell on 2011-10-22 at 14:51 said:}

what do you folks think about the catholic prophecy relating to the names and number of popes, which would indicate that the present Pope Benedict is either the last or penultimate?

To perennial, I agree sensationalism is unhelpful and that truth is often veiled by the reporting of history, but in the present context it's unhelpful to be in denial. The Church is, I believe, ready and willing to atone for its sins, and would make more progress if all her members could admit the faults made. The Lord is merciful, the church has to truth in Him as much as any of us must in contemplating our personal salvation.

As for bringing it into line with tradition, how do you think this idea might relate to the `magical circle' being transformed into a mystical spiral?

\hfill

\texttt{Charlotte Cowell on 2011-10-22 at 14:52 said:}

trust in Him, I meant. To adopt the stance that `we' (ie, the church) have been without sin while all others gathered spots and stains without end is unreasonable (and untrue)

\hfill

\texttt{HOO on 2011-10-22 at 15:05 said:}

Vatican III: The Return of the Jedi!

\hfill

\texttt{Cologero on 2011-10-22 at 16:57 said:}

Something like that. To be more precise, the return of the Mage. I'm thinking of a type of Aghori, although in a higher direction, yet still in revolt against the modern world. For example, \textit{Antisocial Behaviour}\footnote{\url{http://www.meditationsonthetarot.com/antisocial-behaviour}}, or some may prefer the Druid point of view: \textit{Lesson in Practical Magic}\footnote{\url{https://thearchdruidreport.blogspot.com/2011/10/lesson-in-practical-magic.html}}.

\hfill

\texttt{Charlotte Cowell on 2011-10-22 at 17:20 said:}

Return of the Jedi/Mage -- I like the sound of that, it would be most welcome!

\hfill

\texttt{Perennial on 2011-10-23 at 02:43 said:}

Charlotte, If by denial you mean to say I am using history to deny fantasy, then yes, yes I am. I am not even a member of the Roman Catholic Church, so I have no stake in it's defense. However, I do not abide by pseudo-histories masquerading as fact. The fact that you defend the possibility of such troubles me. It is the teaching of the Roman Church that the Church is a perfect society constituted by God (Immortale Dei, Pope Leo XIII) and therefore cannot sin. A collective, in fact, cannot sin. Otherwise, we should expect you to apologize for slavery and the Boer War. The Church therefore has nothing to atone for. I do not see other religious groups ``atoning'' for anything, except for the odd apology of the Amish regarding the Holocaust, nor should they. One should not apologize for things one did not do nor on behalf people. I disagree with this line of reasoning. I cannot apologize for your actions anymore then you can apologize for mine. Certainly, we must trust in His grace, but that is only by committing to positive action. Believing the Church has some backlogue of apologies elicites only sentimentality, and little more. In addition, these non-crimes the Church must ``atone'' for only help the enemies of Tradition, and therefore are even less worthy of discussion. Although I know your heart means well, Charlotte, I take strong, strong difference with the position you are taking, which was indentical to Paul VI and John Paul II. It is not the direction the Church should take in order to make that positive step forward.

\hfill

\texttt{Perennial on 2011-10-23 at 02:47 said:}

Cologero, perhaps what we need is a return to innocence.

\hfill

\texttt{Charlotte Cowell on 2011-10-23 at 07:04 said:}

well I do apologise for slavery and the Boer War and every other evil perpetrated by mankind... .aren't we supposed to pray for each other as well as ourselves? The collective power of prayer, as you know, is immeasurably strong if properly focused. And as for other religious groups (especially Christian), again and again on here it's pointed out that they're less perfect than the catholic church, so what more can we expect? But if the catholic church is true custodian of God's word on Earth, then the bar is set pretty high. Christ committed no sins whatsoever yet not only did he pray for us, he died for us. The church is charged with using Him as Her model for behaviour. Having said all that, I do take your point, which is that you believe the catholic church has done very little wrong in the past. the trouble as far as the world is concerned (which in this context is the salient point) is that this is a matter of opinion and very many people are enraged by what they perceive as wrongdoings (mostly past) on the part of the church. I don't think lambs should be thrown to mad dogs or pearls after swine, but there is a lot to be said for always standing on the moral high ground. (Actually I think Benedict is trying to reclaim this).

Your next comment interests me -- actually it gives me hope, although I doubt this was your intention -- with respect to Paul VI and John Paul II. I presume this is because they were `modernisers'? (If not, then apologies if I now shoot off on the wrong track). As I've mentioned here before, I've experienced a terrible crisis of faith (in a catholic sense, which is the faith I was `converted' into) as I learned about the changes made amid great resistance to the Church in the 60s. I don't ask you to plough through it all again here as I know you've `done it to death', but suffice it to say that the melting pot of worms opened up by this controversy leads into pretty much every black conspiracy known to mankind -- actual or imaginary -- the real heart of darkness in our world. I used to be of a very nervous disposition and when I discovered the extent to which John Paul II -- who I literally adored without reservation -- was considered an antipope almost on a par with anti-Christ by some, it definitely gave me the fear. Not at first, but later in the 90s following all the child abuse allegations. I've no idea why they changed the mass, it can only seem like a bad idea and I don't understand it.

On the other hand, I thought of this last night with respect to the literal interpretation of universalism and the embracing of all other faiths that JPII in particular was (in)famous for. Firstly I will say that he managed to reach me at this point in time -- I was (to all practical intents and purposes) a pagan witch, although I drew the line at robes and pentacles because I certainly didn't have the patience for ceremonial at 18. So there you have one of the positive consequences of casting the net so widely, although it also should be said that he did not accomplish this alone, but via an `outreach' worker -- a Benedictine novice monk. This latter was far more instrumental than the former in my conversion, but nonetheless it occurred during JPII's ministry and I don't tend to see accidents in such things, I see a wondrous design. (Although there CAN be accidents, I realise).

Now, however, the world is faced with a spiritual problem of monumental dimensions as we move towards the next phase of human evolution. Which branch of the Y will we -- as a race -- choose: the left-hand path to grey oblivion on dead planets surrounded by gadgets, a race divided neatly into a mass of morons governed by ammoral elites, or the right road to the golden age that we all know is possible. There waiting for us in fact?

Well, we all know what we want and why, but how to achieve this. We're in a situation where there haven't just been wars on earth there have been terrible wars in heaven. The Holy of Holies has been desecrated, the House of God has been stormed by legions of Satan, and everyone on earth has been required to take a stand on one side or another, consciously or otherwise. What act of divine magic do you think is required of the hopeful citizens of New Jerusalem, followers of Michael? For many it is a matter of a passive end to wrong doing and total embracing of what's good and right. (This stands true for all of us as humans).

But do you think initiates, priests and warriors in God's army get off so lightly; to what end are magical powers to be directed? Who amongst us will take on the dark and bloodthirsty form of Quetzalcoatl so the plumed serpent settles on the right side? Who will observe that one of the world's largest and most powerful group of initiates (Tibetan Buddhists) is utterly divided (certainly astrally and etherically since WWII) into workers of light and those who've been converted to the dark side, and take steps to bring round the renegades? Who will take a compassionate stance before the followers of Hoodoo -- reduced to the status of animals by neighbours and foreigners alike -- in order to harness to great effect the magical elementals currently screaming in agony as `Orthodox Christians' set the world on fire?

What does an alchemist do but take dark materials and transform them into spiritual gold? If the human race is to evolve as it should, towards a higher spiritual plane, the alchemical work must go into overdrive without reservation. This is a fact or everyone will sink into polluted Earth where freedom is nothing more than a memory. It may not even be a memory because the majority of people will be so utterly brainwashed that the collective unconscious will be shrouded forever by chemical clouds. This is why the church and all its children must be unified, no matter what's gone on in the past -- it's a test of the powers of forgiveness, but we're all aware that in fact at the end of it all, after seemingly endless trials and lessons learned on the pillar of severity, genuine atonement brings us over to the side of mercy. In this respect I care much less who and what I'm apologising for, because I'm completley and utterly focused on the transformation. Many Christians believe that the promised land won't come until each and every individual is ready for it, so it's a practical (not even altruistic) matter of hurrying things along as quickly as possible so we can all be bathed in eternal love and light once again!

But I don't know why they changed the mass, I'm working on that one. As for whether other religions also need rise to the occasion, unequivocally yes. I'd be more than happy to stand in front of the chief Rabbi of Jerusalem and say `It's about time' -- I've done it plenty of times and you know what, plenty of the devout Jews I spoke to about this were in total agreement that the chosen people need to share their truth and lead the world in the way they should, which is spiritual. As for the Muslims, the poor Sufis and other purely spiritual sects are having a rough time of it, something needs to make them get their acts together. If only handing over Jerusalem would make them see it might be worth it, but I fear that wouldn't be the case without divine intervention.

\hfill

\texttt{Charlotte Cowell on 2011-10-23 at 07:17 said:}

Perennial, I hate to wave a new age flag right at you but there IS this with respect to Innocence:

\url{https://www.youtube.com/watch?v=UYOVbUygRls}

and this :-)

\url{https://www.youtube.com/watch?v=NTozCB3cueA&feature=related}

\hfill

\texttt{Perennial on 2011-10-24 at 01:08 said:}

Actually, Charlotte, this is the video I was referring too! How did you know? :-) I think the lyrics say it all, do they not?

\hfill

\texttt{Perennial on 2011-10-24 at 01:35 said:}

The Return to Innocence one, of course.

\hfill

\texttt{Charlotte Cowell on 2011-10-24 at 11:07 said:}

Perennial, I am an unabashed lover of Enigma, I think most of their lyrics say it all -- re the `Stargate', try Morphing Thru Time, it works wonders!

I first heard their music during the Trinity term of my second year at university (1996). It was early morning and I was with my Benedictine friend, who was extremely musical and in fact helped me to `hear' (rather than just feel) music for the first time. Hear the words that is.

Anyways, after being in silence for quite some time, out of the blue I asked him... .'Where do we come from?' (the question emerged out of nowhere). Then lo and behold, a second later, `We came out from the deep... .' resounded in full stereo by way of response.

I've loved Enigma ever since -- New Age or no, I do think their music is magical!

\hfill

\texttt{Sparrow on 2015-11-09 at 06:54 said:}

It appears that there are many White Nationalists who follow Gobineau's line of thinking that caste equals biological race. I've never understood how they explain the fact there exists homogeneous Traditional civilizations around the world who have had racially pure (from a biological perspective) lower castes such as Britain and Japan.

Their worldview essentially means that caste means nothing. A shudra who is born with fair hair can become a kshatriya because of biology. The Nazis embodied this philosophy by emphasizing the purity of the peasants.


\hfill

\end{sffamily}
\end{footnotesize}

